\begin{codebox}
\codeline{\textcolor{cmtgray}{\#\ ==============================}}
\codeline{\textcolor{cmtgray}{\#\ 4.\ LLM\ 与本体的互补结构}}
\codeline{\textcolor{cmtgray}{\#\ ==============================}}
\codeline{\textcolor{cmtgray}{\#\ \ \ LLM\ 擅长：自然语言理解、意图识别、文本生成、模糊匹配}}
\codeline{\textcolor{cmtgray}{\#\ \ \ 本体擅长：事实存储、一致性检验、逻辑推理、可解释溯源}}
\codeblank{}
\codeline{\textcolor{cmtgray}{\#\ \ \ 分工模式：}}
\codeline{\textcolor{cmtgray}{\#\ \ \ 用户自然语言\ \(\rightarrow\)\ [LLM]\ 解析意图、生成查询}}
\codeline{\textcolor{cmtgray}{\#\ \ \ \ \ \ \ \ \ \ \ \ \ \ \ \(\rightarrow\)\ [本体/知识图谱]\ 检索事实、校验约束、执行推理}}
\codeline{\textcolor{cmtgray}{\#\ \ \ \ \ \ \ \ \ \ \ \ \ \ \ \(\rightarrow\)\ [LLM]\ 把结构化结果组织成自然语言回答}}
\codeline{\textcolor{cmtgray}{\#\ \ \ （实现模式详见第8章：GraphRAG、Text2SPARQL、本体引导Agent）}}
\codeblank{}
\end{codebox}
